\section{Minimax risk for the budget-value curve}

The shadow-price representation in \cref{obj:prop:dual-representation} reduces estimation of the budget-value curve to estimation of one process. The resulting minimax risk is matched for every interval \(I=[b_0,1-b_0]\), including the full budget range and the single half-budget value.

\begin{theoremv}{thm:matched-curve-frontier}[Matched curve minimax frontier]*
Fix an overlap parameter \(\epsilon\in(0,1/2)\). There exist constants \(0<c_\epsilon\le C_\epsilon<\infty\) such that, for every \(n\ge1\), \(d\ge2\), and \(b_0\in[0,1/2]\), the curve minimax risk of \cref{obj:def:curve-risk} satisfies
\[
c_\epsilon\min\left\{1,\frac{d}{n\log(ed)}\right\}
\le \mathfrak R_{n,d,\epsilon,b_0}^{\mathrm{curve}}
\le C_\epsilon\min\left\{1,\frac{d}{n\log(ed)}\right\}.
\]
For every measurable scalar estimator \(\widehat V_{1/2}\) based on the \(n\) observed records, there is a full-data law \(P\) with the following properties:
\begin{itemize}
\item \textbf{(Causal model.)} \(P\in\mathcal M_{d,\epsilon}\), and each occupied cell has \(p_j>0\), \(\tau_j\in[1/8,3/8]\), and \(\tau_j>0\).
\item \textbf{(Values and binding capacity.)} \(V_1(P)=1/2\), and there is a policy \(\pi\in[0,1]^d\) such that
\[
\sum_{j=1}^{d}p_j\pi_j=\frac12,
\qquad
V_{1/2}(P)=B(P)+\sum_{j=1}^{d}p_j\pi_j\tau_j.
\]
\item \textbf{(Scalar risk.)} Under \(n\) independent observations from the observed margin of \(P\),
\[
\mathbb E_P\!\left[\bigl(\widehat V_{1/2}-V_{1/2}(P)\bigr)^2\right]
\ge c_\epsilon\min\left\{1,\frac{d}{n\log(ed)}\right\}.
\]
\end{itemize}
When \(b_0=0\), every \(P\in\mathcal M_{d,\epsilon}\) also satisfies the full-budget identity
\[
V_1(P)=B(P)+\sum_{j=1}^{d}p_j\max\{\tau_j,0\},
\]
where \(V_b(P)\) is the budget value of \cref{obj:def:budget-value}.
\end{theoremv}
% CAUSALSMITH-CITED-SCOPE-BEGIN thm:matched-curve-frontier
\verificationfootnotetext{\textbf{Formalization scope.} This result uses the published conclusion \citep{JiaoHanWeissman2018} (Theorem 3 and the unknown-both-distributions reduction following Theorem 6). The cited source proof is not formalized here: Lean verifies its use and all remaining steps. If that cited conclusion is false, the portions of this result that depend on it are not certified.}
% CAUSALSMITH-CITED-SCOPE-END thm:matched-curve-frontier

The upper bound applies to all budgets in the chosen interval simultaneously. In the paired family, the half-budget coordinate attains the stated risk order with positive occupied-cell effects, binding capacity, and a full-capacity value fixed at \(1/2\). This family makes the allocation problem visible through the ranking of cell benefits. Simultaneous curve estimation therefore attains the unrestricted full-capacity scalar minimax order under the stated conditions. The reduction uses the bound for estimating the \(\ell_1\) distance between two unknown discrete distributions from \citet[Theorem~3 and the unknown-both-distributions reduction following Theorem~6]{JiaoHanWeissman2018}.

The proof has three steps. First, the shadow-price identity reduces all budget values to one threshold process, and taking a minimum is a contraction in supremum norm. Second, the estimator below controls that process cell by cell: empirical masses suffice for a small alphabet, while pilot localization and factorial moments handle a growing alphabet. Third, the paired family converts the half-budget value into a discrete two-sample \(L_1\) functional, giving the matching lower bound. The appendix contains the polynomial and process estimates behind the second step.

The exact estimator and its uniform upper-bound theorem are stated in the technical appendix, \cref{obj:def:jf-estimator,obj:thm:curve-upper}. Its statistical construction is summarized here.



The estimator first estimates one threshold process across shadow prices, then takes a shadow-price minimum to obtain each budget value. For a small alphabet it inserts empirical cell masses into that process; with a larger alphabet and adequate sample size, pilot counts localize each mass and evaluation counts estimate a local Jackson polynomial by factorial moments. Clipping controls the cell estimates before they are summed, and the sparse-sample branch reports a constant curve. The displayed construction uses an exact conditional average over auxiliary permutations and markings. Its polynomial degree and pilot-width constants serve the risk proof: at ordinary alphabet sizes the degree remains two and the pilot intervals are broad. The rate is a theoretical benchmark for this exact estimator.

In the Jackson--factorial branch, write \(L=\log(ed)\). Pilot counts set cell-specific localization widths; on each resulting rectangle, the Jackson approximation error is bounded by radius-dependent terms divided by the polynomial order \(K\) and \(K^2\) \cref{obj:lem:jackson-boundary-adaptive}. Independent evaluation counts turn the polynomial terms into factorial-moment estimates. The resulting cell bias is bounded by \(C_\epsilon\{\sqrt{p_j/(nL)}+1/(nL)\}\) uniformly over shadow prices, while the centered good-pilot process has expected squared supremum error of order \(d/(nL)\); the bad-pilot envelope has an exponentially small factor in \(L\) \cref{obj:thm:integrated-process-certificate}. Summing the bias over cells and combining it with these fluctuation bounds yields the same process-risk order. The empirical branch covers small alphabets, and the constant branch covers the sparse-sample regime.



Conditional averaging makes the exact curve-valued estimator a function of the observed sample. By \cref{obj:prop:dual-representation}, the error in each budget value is bounded by the uniform error of the estimated threshold process. The same process estimate therefore controls the supremum over \(I\).



The bound measures the squared largest error across the interval before taking expectation, so it also applies when a budget is selected after viewing the estimated curve. The process bound behind it is given in the appendix; the dual identity transfers that bound to budget values without an additional cost for considering several budgets.

The matching converse can be stated directly for estimation at half capacity. It retains positive effects throughout the occupied cells and fixes the value under full provision.

\begin{theoremv}{thm:capacity-active-lower}[Capacity-active risk lower bound]*
Fix an overlap parameter \(0<\epsilon<1/2\). There exists \(c_\epsilon>0\) such that, for every \(n\geq1\), \(d\geq2\), every collection of observed-data laws satisfying \cref{obj:ass:iid-sampling} for each full-data law, and every measurable scalar estimator \(T\) based on the \(n\) observed records, there is a full-data law \(P\) satisfying:
\begin{itemize}
\item \textbf{(Model class.)} \(P\in\mathcal M_{d,\epsilon}\).
\item \textbf{(Full-budget value.)} \(V_1(P)=1/2\), with \(V_b(P)\) as in \cref{obj:def:budget-value}.
\item \textbf{(Effects in occupied cells.)} For every cell \(j\), \(p_j>0\) implies \(1/8\leq\tau_j\leq3/8\).
\item \textbf{(Binding capacity.)} The half-budget constraint is binding for \(P\).
\end{itemize}
For this law,
\[
E_P\!\left[(T-V_{1/2}(P))^2\right]
\geq c_\epsilon\min\left\{1,\frac{d}{n\log(ed)}\right\}.
\]
\end{theoremv}
% CAUSALSMITH-CITED-SCOPE-BEGIN thm:capacity-active-lower
\verificationfootnotetext{\textbf{Formalization scope.} This result uses the published conclusion \citep{JiaoHanWeissman2018} (Theorem 3 and the unknown-both-distributions reduction following Theorem 6). The cited source proof is not formalized here: Lean verifies its use and all remaining steps. If that cited conclusion is false, the portions of this result that depend on it are not certified.}
% CAUSALSMITH-CITED-SCOPE-END thm:capacity-active-lower

Positive treatment effects make the half-budget allocation use its available capacity, while differences across cells determine which recipients are most valuable. The theorem establishes the lower risk benchmark in this capacity-active setting even with the full-budget value fixed. Its reduction again invokes the discrete \(\ell_1\) bound of \citet[Theorem~3 and the unknown-both-distributions reduction following Theorem~6]{JiaoHanWeissman2018}.

Two special cases put the general rate and the budget curve in familiar terms: a fixed covariate alphabet and a common positive effect across occupied cells.

\begin{propositionv}{prop:fixed-alphabet-reduction}[Fixed alphabet rate and identity]*
Fix the overlap parameter \(\epsilon\) and the known covariate alphabet size \(d\) under the following conditions:
\begin{itemize}
\item \textbf{(Overlap.)} \(0<\epsilon<1/2\).
\item \textbf{(Alphabet.)} \(d\ge 2\).
\end{itemize}
For the curve minimax risk in \cref{obj:def:curve-risk}, there are constants \(c,C\) with \(0<c\le C\) such that, for every \(n\ge1\) and every \(b_0\in[0,1/2]\),
\[
\frac{c}{n}\le
\mathfrak R_{n,d,\epsilon,b_0}^{\mathrm{curve}}
\le\frac{C}{n}.
\]

Moreover, let \(P\) belong to the causal model class \(\mathcal M_{d,\epsilon}\) in \cref{obj:def:model-class}. Write \(p_j\) for the mass of cell \(j\), \(\mu_{aj}\) for its arm-\(a\) potential-outcome mean, and \(\tau_j=\mu_{1j}-\mu_{0j}\) for its treatment effect. If some \(c>0\) satisfies \(\tau_j=c\) in every cell with \(p_j>0\), then, for every \(b_0\in[0,1/2]\) and every \(b\in[b_0,1-b_0]\), the budget value in \cref{obj:def:budget-value} satisfies
\[
V_b(P)=B(P)+cb,
\qquad
B(P)=\sum_{j=1}^{d}p_j\mu_{0j}.
\]
\end{propositionv}
% CAUSALSMITH-CITED-SCOPE-BEGIN prop:fixed-alphabet-reduction
\verificationfootnotetext{\textbf{Formalization scope.} This result uses the published conclusion \citep{JiaoHanWeissman2018} (Theorem 3 and the unknown-both-distributions reduction following Theorem 6). The cited source proof is not formalized here: Lean verifies its use and all remaining steps. If that cited conclusion is false, the portions of this result that depend on it are not certified.}
% CAUSALSMITH-CITED-SCOPE-END prop:fixed-alphabet-reduction

With the alphabet and overlap parameter fixed, the matched rate is \(1/n\). Under a common positive effect, every unit of treatment capacity adds the same population value, giving a linear budget-value curve.
In the common-effect clause of \cref{obj:prop:fixed-alphabet-reduction}, the symbol \(c\) denotes an arbitrary positive treatment effect; it is independent of the lower comparison constant also written \(c\) in the risk bound.
